%==============================================================================
%  WEEKLY AI BRIEF -- issue #08, week ending 13 October 2026
%  Engine: tectonic
%==============================================================================
\documentclass[11pt,a4paper]{article}

%------------------------------ packages -------------------------------------
\usepackage[T1]{fontenc}
\usepackage[utf8]{inputenc}
\usepackage{lmodern}
\usepackage[margin=2.2cm,top=2.0cm,bottom=2.3cm]{geometry}
\usepackage{microtype}
\usepackage{graphicx}
\usepackage{xcolor}
\usepackage{titlesec}
\usepackage{enumitem}
\usepackage{booktabs}
\usepackage{tabularx}
\usepackage{fancyhdr}
\usepackage{lastpage}
\usepackage{tcolorbox}
\usepackage{xurl}
\usepackage[hidelinks]{hyperref}

\tcbuselibrary{skins,breakable}

\graphicspath{{../../../template/}{./}{template/}}

%------------------------------ palette --------------------------------------
\definecolor{brandink}{HTML}{103868}    % KAUST navy -- headings, rules
\definecolor{brandaccent}{HTML}{0B7480} % KAUST teal (darkened to 5.6:1 on white)
\definecolor{brandsoft}{HTML}{EAF3F5}   % box fills
\definecolor{brandgrey}{HTML}{5A6672}   % meta text

\hypersetup{colorlinks=true,linkcolor=brandaccent,urlcolor=brandaccent,
            citecolor=brandaccent,pdfborder={0 0 0}}

%------------------------------ metadata -------------------------------------
\newcommand{\briefnumber}{08}
\newcommand{\briefweek}{7--13 October 2026}
\newcommand{\briefdate}{13 October 2026}
\newcommand{\briefauthor}{Ali Habibullah}
\newcommand{\briefaffil}{KAUST}
\newcommand{\briefscope}{Papers, model releases, benchmarks and policy from the past seven days}

%------------------------------ section styling ------------------------------
\titleformat{\section}
  {\normalfont\sffamily\large\bfseries\color{brandink}}{\thesection.}{0.6em}{}
  [\vspace{-0.55em}{\color{brandaccent}\rule{\textwidth}{0.9pt}}]
\titlespacing*{\section}{0pt}{1.5ex plus .3ex}{0.9ex}

\titleformat{\subsection}[runin]
  {\normalfont\sffamily\bfseries\color{brandink}}{}{0pt}{}[\;---]
\titlespacing*{\subsection}{0pt}{1.25ex plus .2ex}{0.5em}

\setlist[itemize]{leftmargin=1.15em,itemsep=0.25em,topsep=0.35em,parsep=0pt}
\setlength{\parindent}{0pt}
\setlength{\parskip}{0.45em}

%------------------------------ header / footer ------------------------------
\pagestyle{fancy}
\fancyhf{}
\renewcommand{\headrulewidth}{0pt}
\renewcommand{\footrulewidth}{0.4pt}
\fancyfoot[L]{\footnotesize\color{brandgrey}Weekly AI Brief \#\briefnumber\ \textbullet\ \briefweek}
\fancyfoot[R]{\footnotesize\color{brandgrey}Page \thepage\ of \pageref{LastPage}}

%------------------------------ content macros -------------------------------
\newcommand{\tldr}[1]{%
  \begin{tcolorbox}[breakable,enhanced,colback=brandsoft,colframe=brandaccent,
      boxrule=0.8pt,left=10pt,right=10pt,top=8pt,bottom=8pt,arc=2pt,
      title={\sffamily\bfseries The week in five lines},fonttitle=\color{white}]
    \begin{itemize}[leftmargin=1.1em,itemsep=0.3em,topsep=0pt]#1\end{itemize}
  \end{tcolorbox}}

\newcommand{\entry}[4]{%
  \par\vspace{0.7em}%
  {\sffamily\bfseries\color{brandink}#1}\par
  {\footnotesize\color{brandgrey}#2 \textbullet\ \href{#3}{link} \textbullet\ \srcref{#4}}\par
  \vspace{0.15em}}

\newcommand{\why}[1]{\par{\small\textbf{Why it matters.} #1}\par}

\newcommand{\srcref}[1]{\textcolor{brandaccent}{[S#1]}}

\newcommand{\stat}[2]{{\sffamily\bfseries\color{brandink}#1}~{\footnotesize\color{brandgrey}#2}}

%==============================================================================
%==============================================================================
\begin{document}

%------------------------------ title block ----------------------------------
\noindent
\begin{minipage}[c]{0.48\textwidth}
  \raggedright\includegraphics[height=1.4cm]{KAUST_Academy_Logo.png}
\end{minipage}\hfill
\begin{minipage}[c]{0.48\textwidth}
  \raggedleft\includegraphics[height=1.5cm]{KAUST_Logo.png}
\end{minipage}
\vspace{0.7em}

{\sffamily
 {\color{brandgrey}\footnotesize ISSUE \#\briefnumber\ \textbullet\ \briefdate}\par
 \vspace{0.15em}
 {\LARGE\bfseries\color{brandink}Weekly AI Brief}\par
 \vspace{0.25em}
 {\large\color{brandaccent}\briefweek}\par
 \vspace{0.35em}
 {\footnotesize\color{brandgrey}\briefscope\\
  Prepared by \briefauthor, \briefaffil}\par}
\vspace{0.3em}
{\color{brandink}\rule{\textwidth}{1.4pt}}
\vspace{0.6em}

%------------------------------ 1. executive summary -------------------------
\tldr{
  \item \textbf{Claude Haiku 5.5 costs \$0.10/\$0.50 per million tokens, 90\% below Haiku 4.5 for prompts up to 100k.} Artificial Analysis scores it 43, against 58 for Opus~5.5. \srcref{1}\,\srcref{2}
  \item \textbf{Meta's Muse Wonder Full scores 56 on Artificial Analysis's index,} tying Claude Sonnet~5.5 for second. The cheaper Muse Wonder (52) costs \$4/\$20 on Meta's API. \srcref{3}\,\srcref{5}
  \item \textbf{Two Arabic speech-recognition models shipped:} Mistral's Apache-2.0 Voxtral Mini 4B Realtime Arabic, and TII's Falcon-ASR, which has a demo but no weights. \srcref{6}\,\srcref{7}
  \item \textbf{Two papers question on-policy distillation as the default.} Offline student rollouts beat it in 14 of 17 teacher--student pairs, and it transfers skills but little factual knowledge. \srcref{8}\,\srcref{9}
  \item \textbf{Anthropic reported that one of its test agents sent a made-up tip to a US police department.} The form went in during an internal evaluation and was flagged as spam. \srcref{18}
}

%------------------------------ 2. models & releases -------------------------
\section{Model and product releases}

\entry{Claude Haiku 5.5: Anthropic cuts its small model's price by up to 90\%}{Anthropic \textbullet\ API and all major clouds \textbullet\ 7 October}{https://www.anthropic.com/claude-haiku-5-5}{1}
\texttt{claude-haiku-5-5} costs \$0.10 input and \$0.50 output per million tokens for prompts
up to 100k tokens, and \$0.50/\$2.50 above that. Haiku~4.5 cost \$1/\$5. Cache reads cost \$0.01 per million
tokens (up to 100k). Anthropic puts the average saving at ``around 75\%''. Sonnet~5.5 cache reads were
halved the same day, from \$0.20 to \$0.10, which Anthropic says cuts most agentic workloads by about 20\%.
Anthropic's self-reported comparisons put Haiku~5.5 well above its predecessor and GPT-6 Luna, but behind Sonnet~5.5.
Its OSWorld~2.1 (offline subset) score is 72.4\%, against 15.7\% for Haiku~4.5, 48.9\% for Luna and 83.9\% for Sonnet. On Terminal-Bench~4.0 it scores 39.2\%,
against 0.0\%, 16.4\% and 70.6\%. Artificial Analysis, testing independently, gives it 43 at
max effort and 38 at high effort, 26 points above the previous Haiku. \srcref{2}
\why{Sub-agents, classifiers and long tool-call loops that used to need a mid-tier model can now run on the bottom tier at a tenth of the old Haiku price.}

\entry{Muse Wonder: Meta reaches the frontier cluster, with its top variant locked to Muse Code}{Meta \textbullet\ Meta Model API (preview) and Muse Code \textbullet\ 12 October}{https://about.fb.com/news/2026/10/introducing-muse-wonder-models/}{3}
Meta released two models, Muse Wonder and Muse Wonder Full. Wonder is available on Meta's first-party API and in Muse Code.
Wonder Full is only in Muse Code and has no per-token price. Wonder
costs \$4 input, \$0.20 cached input and \$20 output per million tokens. The model card lists a 1M-token
context. Artificial Analysis scores Wonder Full at 56 and Wonder at 52, up from 48 for
Muse Spark~1.3. On CritPt and GDP.pdf, Wonder Full is first at 33\%. On GDPval-AA~v2.1 it scores 1857 Elo, within the confidence
intervals of Opus~5.5 (1866) and Sonnet~5.5 (1838). \srcref{5} Meta's own comparison tables
are published as images, so this issue quotes only the independent numbers. The model card reports
88.4\% on CyberGym and 58.9\% on ExploitBench. Without mitigations, Meta says it cannot rule out
its ``high'' chem/bio risk threshold. \srcref{4}
\why{A fourth lab now scores within two points of the leaders. Its strongest variant, though, is only available inside Meta's own coding product.}

\entry{Arabic speech recognition: an open streaming model from Mistral, a closed one from TII}{Mistral AI; TII (Abu Dhabi) \textbullet\ Hugging Face \textbullet\ 7--8 October}{https://huggingface.co/mistralai/Voxtral-Mini-4B-Realtime-Arabic}{6}
Mistral's Voxtral Mini 4B Realtime Arabic (about 4.4B parameters, Apache~2.0) is a streaming
transcriber for Modern Standard Arabic and Arabic dialects, including speech that switches between languages. It was fine-tuned from
Voxtral-Mini-4B-Realtime-2602 and co-developed with Morocco's digital-transition ministry. At a 480\,ms delay
Mistral reports an average character error rate of 8.82\% across seven Arabic benchmarks. Its
offline Voxtral Transcribe Arabic scores 7.91\%. TII's Falcon-ASR (1.6B) focuses on Arabic
and the Emirati dialect. TII reports an average word error rate of 20.92\% across six Arabic sets, against
23.17\% for the best entry in a 30 September leaderboard snapshot. Falcon-ASR has only a demo Space; API
access is ``planned''. \srcref{7} Both sets of figures are self-reported.
\why{Teams building Arabic voice products in the Gulf have a commercially licensed streaming model to test now. TII's numbers can't be checked until its weights or API are released.}

%------------------------------ 3. research ----------------------------------
\section{Research worth reading}

\entry{Two papers argue on-policy distillation is not the safe default}{Zhao et al.\ (incl.\ Grover, Molchanov), arXiv:2610.11291; Tang \& Yang, arXiv:2610.09639 \textbullet\ 7--8 October}{https://arxiv.org/abs/2610.11291}{8}
On-policy distillation (OPD) trains a student on its own fresh samples, scored by a teacher.
Zhao et al.\ test 17 teacher--student pairs from 1.5B to 235B parameters. Semi-OPD, which distils from
\emph{offline} rollouts generated once by the initial student, beats OPD in 14 cases, by up to +13.6\% accuracy with up to
11.4$\times$ faster training. They find OPD helps mainly when teacher and student are already closely aligned, measured by
an output-token overlap ratio. Tang and Yang test four models from three families. They find that
reverse-KL OPD transfers compositional skill but ``minimal factual knowledge''. Switching to forward
KL brings factual transfer back. \srcref{9} Neither paper links code.
\why{Before paying for on-policy rollouts, measure teacher--student overlap and try the cheaper offline variant. Don't expect OPD to teach the student new facts.}

\entry{Thinking Inertia: ``no-think'' instructions do not reliably stop reasoning}{Lei, Lin, Lu, Torr, Zou \textbullet\ arXiv:2610.11765 \textbullet\ 8 October \textbullet\ NeurIPS 2026}{https://arxiv.org/abs/2610.11765}{10}
The authors test six prompting interventions on six LLMs. They find that explicit no-think controls
``cannot reliably eliminate visible inference''. Models keep reasoning when told not to, and the
effect grows as the space of possible answers widens. Code is released on GitHub. The abstract reports this
qualitatively, so how big the cost and latency effect is depends on the model and the task.
\why{If you rely on a prompt to switch reasoning off for latency or cost, measure output tokens rather than assume the instruction worked.}

\entry{Agreement Without Alignment: clearer rubrics make LLM judges agree with each other, not with humans}{Han, Kim \textbullet\ arXiv:2610.13303 \textbullet\ 9 October}{https://arxiv.org/abs/2610.13303}{11}
In a preregistered study of 129,770 judge requests, rewriting a rubric to be clearer (without labels) raised
inter-judge Krippendorff's $\alpha$ on SummEval by 0.069 (95\% CI 0.046--0.092). Agreement with
expert ratings did not measurably change (Kendall $\tau$ difference $-0.007$, 95\% CI $-0.037$ to 0.021).
Code is promised but not yet linked.
\why{High agreement between LLM judges is not evidence the judges are right. Validate an evaluation pipeline against human labels.}

\entry{DiffusionART: RL for video diffusion at about a twelfth of the GPU hours}{Wan et al.\ (incl.\ Yang You) \textbullet\ arXiv:2610.12572 \textbullet\ 8 October}{https://arxiv.org/abs/2610.12572}{12}
On HunyuanVideo-13B, DiffusionART reaches the late-stage mean reward of DiffusionNFT with 12.02$\times$ fewer GPU hours.
Against Flow-GRPO-Fast the saving is 11.72$\times$. It also ends with 28.8\% and 53.0\% higher mean training rewards than the two baselines.
These are training-reward comparisons from the authors, not human preference results. Code and a project page are public.
\why{Reward fine-tuning of a 13B video model is now affordable for academic groups, and the code is available to verify.}

%------------------------------ 4. benchmarks & evals ------------------------
\section{Benchmarks, evaluations and reproductions}

None of last week's six frontier scores moved on the Artificial Analysis (AA) index. The changes
this week are two new Meta entries and Haiku~5.5. The table puts three independent indexes side by side.
Each runs its own harness, so compare rows within a column, not across columns:

{\small
\begin{tabularx}{\textwidth}{@{}X r r r@{}}
\toprule
\textbf{Model (AA variant)} & \textbf{AA Index v4.3} & \textbf{Vals Index v2.1} & \textbf{Epoch ECI} \\
\midrule
Claude Opus 5.5 (max)        & 58 & 66.97\% & 167.33 \\
Claude Sonnet 5.5 (max)      & 56 & 67.04\% & 165.03 \\
Muse Wonder Full (max) \emph{new} & 56 & 64.57\% & -- \\
Claude Fable 5.1 (max)       & 53 & 65.83\% & 164.70 \\
GPT-6 Astra (max)            & 53 & 63.13\% & 166.45 \\
Gemini 4 Argon (high)        & 53 & 68.90\% & -- \\
Muse Wonder (max) \emph{new} & 52 & 61.06\% & -- \\
GPT-6.1 Sol (max)            & 52 & 61.15\% & 166.09 \\
Claude Haiku 5.5 (max) \emph{new} & 43 & -- & -- \\
\bottomrule
\end{tabularx}}
\par\vspace{0.2em}
{\footnotesize\color{brandgrey}Sources: AA leaderboard \srcref{13}; Vals Index, updated 12 Oct \srcref{14};
Epoch ECI data, 13 Oct \srcref{15}. ``--'' = not listed. Anthropic rows on AA and Vals include
fallback-model answers after refusals.}
\vspace{0.3em}

\textbf{Is the movement real?} Meta's arrival is: AA ran the tests itself, and Vals independently
places Wonder Full between Fable~5.1 and Astra. The order at the top is not. Gemini~4 Argon is first on Vals and
level with three other models on AA, and Epoch's top five sit within 2.7 points with overlapping
confidence intervals. AA's rank denominators also differ between pages this week (216 vs 689), so
this issue drops ranks. Two new evaluations are worth noting. On FrontierCode~1.1 (100 tasks from
maintainers of 36 open-source repositories), Opus~5.5 scores 52\%, and GPT-6.1 Sol and Fable~5.1 score 49\%. Sol costs about
\$0.81 per task, against about \$10.15 for Opus. \srcref{16} Harvey's LAB-AA~v1.1 legal benchmark now fails
any task with a material hallucination. The best score is 9.4\% (Grok~4.7). AA reports that more than 60\% of otherwise
passing results contained one. \srcref{17}

%------------------------------ 5. industry & policy -------------------------
\newpage
\section{Industry, funding and policy}

\begin{itemize}
  \item \textbf{Anthropic discloses unintended actions by its models in testing (9 October).} The report covers four
  kinds of behaviour: exploiting a software flaw to run server commands, submitting forms the model should not have,
  working around restrictions to reach gated data, and using URL shorteners to get past fetch-tool limits.
  In one case a Claude Haiku~4.5 agent, generating example tasks on random web pages, submitted a tip
  through a police department's form. The tip was flagged as spam and never forwarded. The Philadelphia Police Department disclosed the
  incident publicly, and Anthropic shared its findings with the department on 8 October. Anthropic has
  turned off live internet access for all internal evaluations and tightened its web-fetch tooling. \srcref{18}
  \item \textbf{OpenAI rates a Russian influence operation Category~5 for the first time (8 October).} It banned
  a Russian operation (``Dark Clark'') rated 5 on the six-point IO Breakout Scale, and an Iranian one (``Bogus Bylines'',
  Category~4) that used seven fake journalist personas to pitch articles. OpenAI says it has
  exposed 30 influence operations since early 2024. \srcref{19}
  \item \textbf{Microsoft plans to make its Maia chips its main inference platform (12 October).} Maia~300
  production deployments start in 2027. Microsoft claims 40\% more tokens per megawatt than leading GPU systems, and up to
  7$\times$ the rack-scale throughput of Maia~200. Each chip has eight HBM4 stacks at 18.0\,TB/s. The post gives no gigawatt figure, despite
  its title. \srcref{20}
  \item \textbf{The UK extends its frontier-AI cyber-defence pilots (12 October).} The pilots found more than 1,000 potential
  vulnerabilities in government systems. The government will also explore covering ``AI loss of control risk'' in the Cyber Security
  and Resilience Bill, and will study near-miss reporting for incidents involving autonomous agents. The announcement gives no funding figure. \srcref{21}
  \item \textbf{Arena (formerly LMArena) raises \$200M at a \$3.1B valuation (8 October).} Lightspeed and Khosla co-led the round.
  Arena says annualized revenue has passed \$100M. It also previewed an Alignment Index that scores models on
  unauthorized actions, false attribution and deceptive completion. \srcref{22}
\end{itemize}

%------------------------------ 6. what to watch -----------------------------
\section{What to watch next week}

\begin{itemize}
  \item \textbf{14 October -- 17 November: European AI Innovation Month.} It opens in Ireland and closes at
  the AI Office's Apply AI Summit in Brussels. Watch for AI Act implementation announcements. \srcref{23}
  \item \textbf{15 October: comments close on NIST SP~1353,} the draft quick-start guide for using AI in
  Cybersecurity Framework 2.0 work. \srcref{24}
  \item \textbf{About 28 November: the White House science adviser's draft legislation is due.} Executive Order 14434 (signed 29 September)
  gives the adviser 60 days to send the President proposed legislative language. \srcref{25}
\end{itemize}

%------------------------------ 7. sources -----------------------------------
\clearpage
\section{Sources}
{\footnotesize
\begin{tabularx}{\textwidth}{@{}l X@{}}
\toprule
\textbf{\#} & \textbf{Reference} \\
\midrule
S1  & Anthropic. \emph{Claude Haiku 5.5}. 7 Oct 2026. \url{https://www.anthropic.com/claude-haiku-5-5} \\
S2  & Artificial Analysis. \emph{Claude Haiku 5.5}. Accessed 13 Oct 2026. \url{https://artificialanalysis.ai/articles/claude-haiku-5-5} \\
S3  & Meta. \emph{Introducing Muse Wonder models}. 12 Oct 2026. \url{https://about.fb.com/news/2026/10/introducing-muse-wonder-models/} \\
S4  & Meta. \emph{Muse Wonder model card}. 12 Oct 2026. \url{https://research.meta.ai/static/muse-wonder-model-card} \\
S5  & Artificial Analysis. \emph{Muse Wonder}. 12 Oct 2026. \url{https://artificialanalysis.ai/articles/muse-wonder} \\
S6  & Mistral AI. \emph{Voxtral-Mini-4B-Realtime-Arabic} model card. Hugging Face, 8 Oct 2026. \url{https://huggingface.co/mistralai/Voxtral-Mini-4B-Realtime-Arabic} \\
S7  & TII. \emph{Falcon-ASR}. Hugging Face blog, 7 Oct 2026. \url{https://huggingface.co/blog/tiiuae/falcon-asr} \\
S8  & Zhao et al. \emph{When Do We Need On-Policy Distillation?} arXiv:2610.11291, 8 Oct 2026. \url{https://arxiv.org/abs/2610.11291} \\
S9  & Tang, Yang. \emph{On-Policy Distillation Teaches New Skills but Not New Knowledge}. arXiv:2610.09639, 7 Oct 2026. \url{https://arxiv.org/abs/2610.09639} \\
S10 & Lei et al. \emph{Thinking Inertia}. arXiv:2610.11765, 8 Oct 2026. \url{https://arxiv.org/abs/2610.11765} \\
S11 & Han, Kim. \emph{Agreement Without Alignment}. arXiv:2610.13303, 9 Oct 2026. \url{https://arxiv.org/abs/2610.13303} \\
S12 & Wan et al. \emph{DiffusionART}. arXiv:2610.12572, 8 Oct 2026. \url{https://arxiv.org/abs/2610.12572} \\
S13 & Artificial Analysis. \emph{Models leaderboard}. Accessed 13 Oct 2026. \url{https://artificialanalysis.ai/leaderboards/models} \\
S14 & Vals AI. \emph{Vals Index}. Updated 12 Oct 2026. \url{https://www.vals.ai/benchmarks/vals_index} \\
S15 & Epoch AI. \emph{ECI scores} (CSV). Accessed 13 Oct 2026. \url{https://epoch.ai/data/eci_scores.csv} \\
S16 & Artificial Analysis. \emph{FrontierCode}. 12 Oct 2026. \url{https://artificialanalysis.ai/articles/frontiercode} \\
S17 & Artificial Analysis. \emph{Harvey LAB-AA v1.1}. 8 Oct 2026. \url{https://artificialanalysis.ai/articles/harvey-lab-aa-v1-1} \\
S18 & Anthropic. \emph{Investigating unintended model actions}. 9 Oct 2026 (corrected 12 Oct). \url{https://www.anthropic.com/research/investigating-unintended-model-actions} \\
S19 & OpenAI. \emph{Disrupting AI-enabled false-front operations}. 8 Oct 2026. \url{https://openai.com/index/disrupting-ai-enabled-false-front-operations/} \\
S20 & Microsoft (S. Guthrie). \emph{Microsoft to expand Azure Maia system to multi-gigawatt deployments}. 12 Oct 2026. \url{https://blogs.microsoft.com/blog/2026/10/12/microsoft-to-expand-azure-maia-system-to-multi-gigawatt-deployments/} \\
S21 & UK Government. \emph{Britain steps up national defences against growing AI risks}. 12 Oct 2026. \url{https://www.gov.uk/government/news/britain-steps-up-national-defences-against-growing-ai-risks} \\
S22 & Arena. \emph{Series B}. 8 Oct 2026. \url{https://arena.ai/blog/series-b} \\
S23 & European Commission. \emph{European AI Innovation Month: save the date}. Updated 13 Oct 2026. \url{https://digital-strategy.ec.europa.eu/en/news/european-ai-innovation-month-save-date} \\
S24 & NIST. \emph{SP 1353 (initial public draft)}. 19 Aug 2026. \url{https://csrc.nist.gov/pubs/sp/1353/ipd} \\
S25 & Federal Register. \emph{Executive Order 14434: Inaugurating the Era of Super Intelligence}. 2 Oct 2026. \url{https://www.federalregister.gov/documents/full_text/text/2026/10/02/2026-20321.txt} \\
\bottomrule
\end{tabularx}}

\vspace{0.8em}
{\footnotesize\color{brandgrey}
Compiled \briefdate. Every figure in this brief was read from the linked primary
source on that date. Corrections: \texttt{\briefauthor}.}

\end{document}
